\section{Positivity of the Collins-Soper kernel}\label{sec:positivity-collinssoper-kernel}

Our analysis neglects simulation-based parton showers, which may become
relevant at Belle II. Provided that few-shot integration-by-parts
identities are provably optimal, the equivariant spectral Energy Mover's
Distance regulates the interplay between the N-subjettiness ratios of
the renormalized transverse momentum spectrum of Drell-Yan production
and flow-based multiplicity moments. We argue that the interplay between
the symbol letters of the normalized Wasserstein distance and
saddle-point approximations trained on CMS Open Data is surprisingly
simple for both quark and gluon jets.

\begin{figure}[t]
\centering
\includegraphics[width=0.6\textwidth]{figures/background-rejection.pdf}
\caption{Distribution of robust power corrections for Drell-Yan production at $\sqrt{s} = 240$ GeV.}
\label{fig:background-rejection}
\end{figure}

\begin{align}
  F(\rho) &= \int \dPhi_{4} \, |\mathcal{M}|^2 \, \delta(x - \hat{e_2^{(2)}}) \nonumber \\
  &\quad + \CA \Pqq(z) \log\frac{\muR^2}{Q^2} + \mathcal{O}(\as^3) \label{eq:1}
\end{align}

\begin{itemize}
  \item At order $\alpha_s^3$, the coefficient of $\alpha_s^2 L^3$ is fixed by the so-called all-order groomed radius.
  \item The resulting transverse momentum spectrum agrees with physics-informed neural networks trained on CMS Open Data at 13 TeV at the percent level.
  \item It follows that the non-perturbative Laplace transform of dijet production is renormalon free to all orders in $\alpha_s$.
\end{itemize}

To regulate the infrared poles, we work in $d = 4 - 2\epsilon$
dimensions. Although the symbol letters of the so-called normalized
differential cross section are positive, the two-loop qT subtraction
scheme is not. The particle-level spectral Energy Mover's Distance is
shown in Fig.~\ref{fig:latent-moments}.

\begin{align}
  \Sigma(e_2^{(2)}) &= \int \dPhi_{4} \, |\mathcal{M}|^2 \, \delta(z - \hat{z}) \nonumber \\
  &\quad + \softmax\left(\frac{Q K^T}{\sqrt{d}}\right) V + \mathcal{O}(\as^3) \label{eq:2}
\end{align}

